\section{Generalisation under decision censoring}
\label{sec:generalization}

Let $\F$ be a class of scores with $\|f\|_\infty\le B$ and $\hat f$ the
minimiser of the empirical worst-case risk with estimated bounds
$[\hat\plo,\hat\pup]$. \Cref{thm:generalization} (\cref{app:generalization-full})
bounds the excess worst-case risk
$\Rbar_\Gamma(\hat f)-\inf_{\F}\Rbar_\Gamma$ by
$4L\,\Rad_n(\F)+2\ell_{\max}(B)\sqrt{\log(2/\delta)/2n}$ plus a nuisance term
$2\kappa(B)\big(\|\hat e-e\|_{L^1}+\tfrac14\|\hat\eta_1-\eta_1\|_{L^1}\big)$,
where $\eta_1=\logit p_1$ and $\kappa(B)=\sup_{|z|\le B}g_\ell(z)$. The range
$B$ is not free: to contain the population optimum \eqref{eq:bayes-act} the
class must reach $|\logit\pup|\le|\logit p_1|+\log\Gamma$, so
$B=B_0+\log\Gamma$, and \cref{cor:loggamma} makes the complexity term
\textbf{linear in $\log\Gamma$}, not in $\Gamma$; \cref{prop:lower-gen} shows
this dependence is necessary. It is the payoff of parameterising the
sensitivity model on the log-odds scale, where $\Gamma$ acts as a translation
and the map from $p_1$ to the bounds is $1$-Lipschitz, rather than on the
probability scale, where it would have Lipschitz constant $\Gamma$.
Empirically (\cref{sec:experiments}) the uniform deviation decays at the
predicted $n^{-1/2}$ rate and its coefficient is far better explained by
$\log\Gamma$ than by $\Gamma$.
